\documentclass[a4paper,11pt]{article}
% Generated by render_rug.js -- do not edit by hand.
% Package set and page geometry follow
% Econometrics_Economics_Thesis_Paper_template/preamble.tex.
\usepackage[T1]{fontenc}
\usepackage[utf8]{inputenc}
\usepackage{geometry}
\usepackage{amsmath}
\usepackage{amssymb}
\usepackage{graphicx}
\usepackage{pdflscape}
\usepackage{booktabs}
\usepackage{array}
\usepackage{tabularx}
\usepackage{longtable}
\usepackage{textcomp}
\usepackage[sort]{natbib}
\usepackage[colorlinks]{hyperref}
\hypersetup{citecolor=blue,linkcolor=blue,urlcolor=blue}
% setspace and placeins are what the template uses for double spacing and
% for holding floats inside their own section. Neither ships with a base
% TinyTeX install, so each has a fallback that keeps the document
% compiling; install them for exact template fidelity with
%   tlmgr update --self && tlmgr install setspace placeins
\IfFileExists{setspace.sty}{\usepackage{setspace}}{%
  \newcommand{\doublespacing}{\linespread{1.667}\selectfont}%
  \newcommand{\singlespacing}{\linespread{1}\selectfont}}
\IfFileExists{placeins.sty}{\usepackage{placeins}}{%
  \newcommand{\FloatBarrier}{}}
% Long unbreakable compounds ("climate-exposed", "intention-to-treat") are
% dense in this prose; let TeX stretch a line rather than overflow.
\emergencystretch=3em
\setlength{\parskip}{6pt}
\graphicspath{{./}}

\geometry{left=2.5cm,textwidth=16.5cm,top=1in,bottom=1in,headheight=15pt}
\doublespacing
\begin{document}
\begin{titlepage}
\linespread{1}
\begin{center}
\vspace*{2cm}
% \includegraphics[width=0.35\textwidth]{logo.png}\\[1cm]
\LARGE\textbf{Deferred Costs: Climate Shocks, Local Economic Capacity, and Health-Market Ledgers}

\vspace{0.6cm}
\Large Ahwaz Akhtar

\vspace{0.8cm}
\large Doctoral Dissertation --- Essay 1

\vspace{0.4cm}
\today

\vspace{0.4cm}
\hrulefill
\vspace{0.4cm}

\end{center}
\begin{abstract}\noindent\footnotesize Climate shocks impose healthcare costs that continue to emerge after the extreme event has passed. How these costs are observed depends overwhelmingly on the data sources we choose to consult. Using a harmonized county and state panel linking month-level climate data to Medicare records, credit-bureau medical debt, ACA premiums, and county income accounts, this paper explores what healthcare costs follow climate shocks in the United States, and which financial sources record these.  I find that these administrative health records measure the morbidity channel directly: as heat shocks are observed, standardized Medicare spending rises by \$112 per beneficiary in the shock year and \$176 the next, and emergency visits increase by 8-9 per 1,000, among the 65-plus and disabled population, with the largest responses arriving the year after exposure. I further find evidence of uneven financial responses to climate shocks: extreme cold raises the credit-bureau medical-debt share by 0.85 to 1.35 percentage points at a one-year lag, and ACA premiums show no coherent local pass-through. Similarly, for drought, the data rule out pass-through of measured morbidity costs to ACA premiums. I find that the unpriced costs strain household budgets: drought lowers county incomes contemporaneously through a nonfarm channel of \$261 to \$414, stable across multiple choices of estimation windows. Overall, this paper contributes to literature on measurement of deferred morbidity costs directly, and identifies which data sources record such costs for different populations.\end{abstract}
\noindent\hrulefill
\end{titlepage}

\section*{Glossary of Abbreviations}
% Glossary of abbreviations for the combined thesis. Static content: every entry
% appears in at least one of the three essays or their exhibits.
\begingroup
\small
\setlength{\tabcolsep}{6pt}
\begin{longtable}{@{}p{2.4cm}p{12.2cm}@{}}
\toprule
\textbf{Abbreviation} & \textbf{Expansion} \\
\midrule
\endfirsthead
\toprule
\textbf{Abbreviation} & \textbf{Expansion} \\
\midrule
\endhead
\bottomrule
\endfoot

ACA   & Affordable Care Act; here, the individual health-insurance marketplace it created \\
ACS   & American Community Survey (US Census Bureau) \\
AQI   & Air Quality Index (US Environmental Protection Agency) \\
ATT   & Average treatment effect on the treated \\
BEA   & Bureau of Economic Analysis \\
BKY   & Benjamini--Krieger--Yekutieli sharpened false-discovery-rate procedure \\
CDC   & Centers for Disease Control and Prevention \\
CDD   & Cooling degree days; the heat-exposure measure \\
CMS   & Centers for Medicare \& Medicaid Services \\
CONUS & Contiguous United States \\
CPI   & Consumer Price Index, used to convert nominal to real dollars \\
CS    & Callaway--Sant'Anna staggered difference-in-differences estimator \\
DiD   & Difference-in-differences \\
DOE   & US Department of Energy \\
DRDID & Doubly robust difference-in-differences estimator \\
ED    & Emergency department \\
ERS   & Economic Research Service (US Department of Agriculture) \\
FE    & Fixed effects \\
HDD   & Heating degree days; the cold-exposure measure \\
HIX   & Health Insurance Exchange; source of marketplace premium data \\
IRS SOI & Internal Revenue Service Statistics of Income, used for migration flows \\
ITT   & Intention to treat \\
LEAD  & Low-Income Energy Affordability Data (US Department of Energy) \\
LP    & Local projection \\
NASHP & National Academy for State Health Policy, source of hospital financial data \\
NHE   & National Health Expenditure Accounts \\
NOAA  & National Oceanic and Atmospheric Administration \\
NPR   & Net patient revenue \\
PDSI  & Palmer Drought Severity Index; extreme drought is a value at or below $-4$ \\
PMPM  & Per member per month, the unit in which premiums are compared to costs \\
pp    & Percentage points \\
PRISM & Parameter-elevation Regressions on Independent Slopes Model, the gridded climate product used for humidity \\
PUF   & Public use file \\
RA    & Rating area, the geographic unit at which marketplace premiums are set \\
RI    & Randomization inference \\
SAHIE & Small Area Health Insurance Estimates (US Census Bureau) \\
SD    & Standard deviation \\
SE    & Standard error \\
SVI   & Social Vulnerability Index (CDC/ATSDR) \\
USDA  & US Department of Agriculture \\
WCB   & Wild cluster bootstrap, with Webb weights \\
\end{longtable}
\endgroup

\clearpage
\FloatBarrier
\section{Introduction}

When a drought or a heat wave ends, the weather normalizes long before the bills arrive. The costs of a climate shock can surface a year or more later in places such as a county's personal income accounts, in Medicare claims, or in collection lines on a credit report. This paper attempts to document which financial outcomes are sensitive to these climate shocks and the extent to which household finances and healthcare expenditure are affected by these shocks.

Delayed incidence is hard to measure for a few reasons. Attribution across years is confounded by factors other than climate; adaptation and migration change the exposed population while the costs are still being realized; and every administrative data source observes a different population under different recording rules. A large literature measures the contemporaneous mortality and output effects of temperature and drought \citep{deschenes2007agricultural,dell2014weather,hsiang2016econometrics}; the financial after-path of a shock through household and health-market data sources is far less mapped \citep{dobkin2018hospital,kluender2021debt}.

I take three steps to create a harmonized panel. First, I measure a morbidity (increased utilization/cost) channel directly in Medicare administrative data. Second, I test whether two financial data sources - credit-bureau medical debt and ACA benchmark premiums - record the increased morbidity-related costs. Third, I document whether the household economic capacity absorbs these additional costs. Climate and pollution exposure come from NOAA's National Centers for Environmental Information (temperature, heating and cooling degree-days, and the Palmer Drought Severity Index, all recorded at the county-climate-division level), from PRISM - Oregon State University's gridded interpolation of weather-station data, used here for dew-point humidity - and from EPA county-level air-quality monitors.

First, I find that administrative health records measure the morbidity channel directly. Extreme heat raises standardized Medicare spending by \$112 per beneficiary in the shock year and \$176 the next - 1.1 and 1.7 percent of the \$10,359 mean - and emergency-department visits by 8-9 per 1,000 beneficiaries, reproducing \citet{deryugina2019mortality} within this panel; cold shocks operate at longer lags, while air quality shocks peak in the year of incidence. The largest responses arrive the year after exposure. The evidence covers the 65-plus and disabled population over 2014--2023, and the hazards that drive it - heat and air quality - reach the Medicare ledger without a farm-income intermediary.

Second, the financial sources record this harm unevenly. Cold shocks raise the credit-bureau medical-debt share by 0.85 to 1.35 percentage points at a one-year lag in the state panel, an association that does not reproduce at the county level, but bureau debt is measurement-fragile: it records hardship only after insurance, billing, and credit-file filters. ACA premiums show no coherent local pass-through at all - coefficients flip sign across the county, rating-area, and state levels, and within-state responses are bounded below roughly 4 to 8 percent of the mean premium. For drought, the data rule out pass-through larger than 51 to 80 percent of the measured morbidity costs. Consistent with the null first stage, 92 to 99 percent of each shock's debt response survives premium adjustment.

The third set of results documents the household capacity behind these ledgers. Drought lowers county incomes in distributed-lag models whose coefficients are stable across estimation windows reaching back to 1990 - between -\$99 and -\$132 per unit of drought severity, with a contemporaneous term of -\$149 (p = <0.01) in the longest window. An event study of the 2012 drought, reported in Appendix A, shows why I do not headline the corresponding event contrast: the raw gap of -\$1,311 survives wild-cluster bootstrap and randomization inference, but roughly \$900 of it reflects farm earnings reverting from their record 2011 commodity-price peak. Event contrasts in agricultural counties conflate the commodity-price cycle with climate damage unless the income ledger is decomposed, and the essay treats that decomposition as part of its contribution.

The paper contributes to three literatures. To health economics, it adds administrative utilization responses measured in the same panel as the economic outcomes, so the health channel is documented by measurement rather than inference \citep[nearest antecedent:][]{deryugina2019mortality}. To household finance and the economics of measurement, it adds the finding that regulated ACA pricing shows no coherent local response to climate-health shocks, and credit-bureau debt records them selectively. Therefore, which source one consults determines whether the cost is visible at all. To the climate-economy literature, it adds a decomposition result: event contrasts in agricultural counties inherit the commodity-price cycle embedded in the income ledger, and separating farm from nonfarm income materially revises what such contrasts measure \citep[nearest antecedents:][]{deschenes2007agricultural,deryugina2017fiscal}.

In this paper, Section 2 describes the institutions and Section 3 the data. Section 4 reports the Medicare morbidity evidence, Section 5 the financial-ledger responses, and Section 6 the household economic capacity results, with the 2012 drought natural experiment and its decomposition reported in full in Appendix A. Section 7 discusses interpretation and limits, and Section 8 concludes.

\FloatBarrier
\section{Institutional and conceptual background}

A simple framework organizes the essay: climate exposure affects households through four margins - health and healthcare utilization, labor productivity and local income, household liquidity, and provider demand - and each margin is observed through a different institution. Medicare records utilization for the 65-plus and disabled population; employer and BEA accounts record employment and income; ACA rating areas record regulated premiums; credit bureaus record unpaid medical bills; and hospitals record uncompensated care. These institutions differ in eligibility, geography, recording rules, and timing, so identical underlying harm can produce different - even opposite-signed - cost responses. Figure~\ref{fig:e1f0} sets out this frame.

\begin{figure}[htbp]
\centering
\caption{Two routes from a climate shock to a financial ledger. The mediated route runs through a health-utilisation or income mechanism. The direct route, drawn heavier, is the one this essay leads with: heat, cold, and air quality reach the Medicare ledger as a measured cost, with no intervening mechanism observed and no farm-income intermediary.}
\label{fig:e1f0}
\includegraphics[width=0.95\linewidth]{../../../Analysis/descriptive/fig_conceptual_model.png}
\end{figure}

Three features of ACA rating matter for our understanding. Premiums are set at the rating-area level by insurers operating a single statewide risk pool per market, subject to state regulatory review. The geographic rating factor may reflect unit costs only - not local morbidity - and federal risk adjustment under Section 153 redistributes claims risk across insurers statewide. Further, the rate calendar is slow: plan year t rates are filed around the middle of year t-1 on claims experience through roughly t-2, and are locked before the plan year begins. Therefore, only lagged shocks can be in the insurer's information set, and any local claims signal is institutionally muted before it can reach a local premium.

Credit-bureau medical debt passes through three filters before it is recorded: an insurance relationship, a billed encounter, and a credit file. Hardship that fails any filter never appears in the ledger. Reporting rules - collection thresholds and reporting delays - also changed over the sample window, shifting the recorded object over time. The third essay shows the recorded response to drought is smaller where uninsurance is higher, the visibility gradient this structure predicts.

Lastly, the 2012 drought was among the most severe and extensive U.S. droughts since at least 1895, and it developed suddenly in late spring, arriving without seasonal warning \citep{hoerling2014drought}. I define extreme drought as a Palmer index at or below -4. The counties experiencing their first panel-era extreme onset in 2012 lie in Georgia, the Mountain West, and the Plains, and their economies range from heavily agricultural to barely agricultural, a range that matters when Section 4 shows the burden operates substantially outside agriculture, and again when Appendix A decomposes the event contrast into farm and nonfarm income.

\FloatBarrier
\section{Data and construction of shocks}

The primary dataset is a county-year panel covering 2011 to 2023, with roughly 3,100 counties entering the estimation sample. A companion state panel extending back to roughly 1996 carries the longer medical-debt and premium series. All dollar values are deflated to 2023 dollars.

Table~\ref{tab:e1t0a} lists the sources that compose the panel along with their temporal coverage. The variables broadly span 4 categories. Climate and environmental exposure come from NOAA's nClimDiv county series, EPA air-quality monitors, and PRISM dew-point grids. Economic capacity comes from BEA county income accounts and the American Community Survey. The health-market ledgers come from Urban Institute credit-bureau records, HIX Compare marketplace premiums, and the CMS Geographic Variation file. The moderators used in the mechanism tests come from ACS industry composition, the DOE LEAD energy-burden data, and USDA ERS farm-dependence measures. Two coverage limits shape what follows. Air-quality monitors operate in 1,194 counties rather than the full panel, so air-quality results describe monitored, disproportionately urban, counties. PRISM grids cover the contiguous states only. The coverage reported in the table is data availability rather than the estimation sample, which is 2011 to 2023 for county outcomes, 2014 to 2023 for Medicare, and the marketplace era for premiums.

\begingroup
\footnotesize
\setlength{\tabcolsep}{4pt}
\sloppy
\begin{longtable}{>{\raggedright\arraybackslash}p{1.8cm} >{\raggedright\arraybackslash}p{4.30cm} >{\raggedright\arraybackslash}p{2.9cm} >{\raggedright\arraybackslash}p{2.9cm} >{\raggedright\arraybackslash}p{3.1cm}}
\caption{Data sources, coverage, and role in Essay 1}
\label{tab:e1t0a}\\
\toprule
Domain & Source & Unit and geography & Coverage & Role in Essay 1 \\
\midrule
\endfirsthead
\multicolumn{5}{l}{\textit{Table \thetable{} (continued)}} \\
\toprule
Domain & Source & Unit and geography & Coverage & Role in Essay 1 \\
\midrule
\endhead
\midrule \multicolumn{5}{r}{\textit{continued on the next page}} \\
\endfoot
\bottomrule
\endlastfoot
Climate & NOAA NCEI (nClimDiv) & County-year; US counties & 1990--2025; 111,852 obs.; 3,107 units & Heat, cold, and drought shock construction \\
Air quality & EPA AQS annual AQI & County-year; Monitored counties & 2008--2025; 18,554 obs.; 1,194 units & Air-quality shock in the Medicare models \\
Humidity & PRISM (Oregon State / NACSE) & County-year; CONUS counties & 2009--2023; 46,620 obs.; 3,108 units & Humidity robustness check only \\
Income & BEA CAINC1 / CAINC5N & County-year; US counties & 1990--2024; 110,215 obs.; 3,149 units & Per-capita income outcome; farm/nonfarm split in Appendix A \\
Employment & Census ACS 5-year (B23025) & County-year; US counties & 2011--2023; 40,154 obs.; 3,100 units & Employment outcome and labor-exposure moderator \\
Medical debt & Urban Institute (credit-bureau, Aug snapshots) & County-year; US counties & 2011--2023; 38,870 obs.; 3,078 units & Debt-ledger outcome; treated as measurement-fragile \\
ACA premiums & HIX Compare individual market & County-year (rating-area price); Marketplace rating areas & 2014--2025; 36,974 obs.; 3,144 units & Pass-through outcome; estimated at the rating-area level \\
Medicare & CMS Geographic Variation PUF & County-year; US counties, 65+/disabled & 2014--2023; 31,429 obs.; 3,156 units & Main result: standardized spending and ED visits \\
Industry mix & Census ACS C24030 & County-year; US counties & 2011--2023; 41,868 obs.; 3,234 units & Climate-exposed non-farm share moderator \\
Energy burden & DOE LEAD 2022 & County (time-invariant); US counties & 3,144 counties & Energy-burden moderator \\
Ag dependence & USDA ERS Typology + BEA CAINC5N & County (time-invariant); US counties & 3,116 counties & Agricultural-channel tests in Appendix A \\
Migration & IRS SOI county-to-county flows & County-year; US counties & 2012--2020; 27,899 obs.; 3,143 units & Selection bound on the drought scar (Appendix A) \\
Health spending & CMS National Health Expenditure Accounts & State-year; US states & 1996--2020; 1,275 obs.; 51 units & Systemic-spending null (contrast motivating household ledgers) \\
Macro controls & BLS / BEA via FRED & State-year; US states & 1996--2025; 1,500 obs.; 50 units & State-panel controls and debt-series context \\
\end{longtable}
\vspace{-0.6em}\begin{minipage}{\textwidth}\footnotesize
Coverage is computed from the delivered files at build time: `Years' is the range over non-missing observations, `Obs.' the count of non-missing observations, and `Units' the number of distinct counties or states with at least one non-missing value. Counts therefore reflect data availability, not the estimation sample, which is 2011--2023 for county outcomes, 2014--2023 for Medicare, and the marketplace era for premiums. County dollar series are in 2023 USD; state dollar series are in 2025 USD.
\end{minipage}
\endgroup


Climate exposure comes from NOAA temperature and degree-day records and the Palmer Drought Severity Index, supplemented by PRISM dew-point humidity for robustness and EPA air-quality data for pollution shocks. I define four binary shock bins against a 1990--2000 climate baseline: extreme heat (High\_CDD) and extreme cold (High\_HDD) mark degree-day totals above the 80th percentile of the baseline-period distribution; extreme drought marks a Palmer index at or below -4; and poor air quality (High\_AQI) marks exceedances of the EPA threshold of 100 AQI. Temperature z-scores beyond $\pm$1.5 against the same baseline define the z-based shocks used in robustness. Continuous precipitation z-scores enter all county models and the entire climate hazard family is documented in Appendix B.

Outcomes are drawn from five sources, each observing a different population under different rules: BEA per-capita personal income; civilian employment; the Urban Institute's credit-bureau medical-debt share; ACA benchmark premiums at the rating-area level (2014--2025); and CMS Geographic Variation files for Medicare spending and utilization among the 65-plus and disabled population (2014--2023). Table~\ref{tab:e1t0b} reports definitions, populations, and observation windows.

\begingroup
\footnotesize
\setlength{\tabcolsep}{4pt}
\sloppy
\begin{longtable}{>{\raggedright\arraybackslash}p{2.7cm} >{\raggedright\arraybackslash}p{5.00cm} >{\raggedright\arraybackslash}p{2.1cm} >{\raggedright\arraybackslash}p{2.7cm} >{\raggedright\arraybackslash}p{2.5cm}}
\caption{Variable definitions, sources, and coverage}
\label{tab:e1t0b}\\
\toprule
Variable & Definition & Source & Coverage & Used in \\
\midrule
\endfirsthead
\multicolumn{5}{l}{\textit{Table \thetable{} (continued)}} \\
\toprule
Variable & Definition & Source & Coverage & Used in \\
\midrule
\endhead
\midrule \multicolumn{5}{r}{\textit{continued on the next page}} \\
\endfoot
\bottomrule
\endlastfoot
\multicolumn{5}{l}{\textit{Climate shocks}} \\
Heat shock & 1 if annual cooling degree days are at or above the national 80th percentile of the 1990--2000 baseline distribution (Binary) & NOAA NCEI & 1990--2026; 116,039 obs. & Main Medicare models; labor-exposure interactions \\
Cold shock & 1 if annual heating degree days are at or above the national 80th percentile of the 1990--2000 baseline distribution (Binary) & NOAA NCEI & 1990--2026; 116,039 obs. & Medicare models; debt ledger \\
Extreme drought & 1 if the annual mean Palmer Drought Severity Index is at or below -4 (Binary) & NOAA NCEI & 1990--2026; 116,039 obs. & Debt ledger; premium pass-through; Appendix A event \\
Air-quality shock & Annual maximum AQI; the shock indicator is Max AQI above 100 (EPA `unhealthy' threshold) (AQI index) & EPA AQS & 2008--2025; 18,554 obs. & Medicare ED-visit models \\
Temperature z-score & Annual mean temperature standardized to each county's own 1990--2000 mean and standard deviation (SD) & NOAA NCEI & 1990--2025; 112,932 obs. & Continuous-dose robustness \\
\addlinespace
\multicolumn{5}{l}{\textit{Outcomes}} \\
Medicare standardized spending & Standardized Medicare payment per beneficiary, price-adjusted by CMS to remove geographic payment-rate variation (USD per beneficiary) & CMS Geographic Variation & 2014--2023; 31,429 obs. & Main result \\
Medicare ED visits & Emergency department visits per 1,000 Medicare beneficiaries (Visits per 1,000) & CMS Geographic Variation & 2014--2023; 31,409 obs. & Main result \\
Per-capita personal income & BEA per-capita personal income, CPI-deflated (2023 USD) & BEA CAINC1 & 1990--2024; 110,215 obs. & Household economic capacity; Appendix A \\
Civilian employment & Count of employed civilians; estimated in logs so effects are proportional rather than county-size driven (Persons (log)) & ACS B23025 & 2011--2023; 40,154 obs. & Employment margin; labor-exposure interactions \\
Medical debt share & Share of credit records with medical debt in collections (Share) & Urban Institute & 2011--2023; 38,870 obs. & Debt ledger (measurement-fragile) \\
Benchmark silver premium & Second-lowest-cost silver marketplace premium, CPI-deflated; aggregated to rating-area by population-weighted mean for estimation (2023 USD per month) & HIX Compare & 2014--2025; 36,974 obs. & Pass-through null and equivalence bounds \\
Lowest bronze premium & Lowest-cost bronze marketplace premium, CPI-deflated (2023 USD per month) & HIX Compare & 2014--2025; 36,974 obs. & Pass-through robustness \\
Health spending per capita & State personal health-care expenditure per capita, CPI-deflated (2025 USD) & CMS NHE & 1996--2020; 1,275 obs. & Systemic-spending null \\
\addlinespace
\multicolumn{5}{l}{\textit{Moderators}} \\
Climate-exposed non-farm share & Baseline share of employment in construction, manufacturing, transportation and warehousing, and utilities; held fixed at its baseline value and entered as a z-score (Share (z-scored)) & ACS C24030 & 3,221 counties & Labor-exposure mechanism \\
Household energy burden & Household-weighted energy expenditure as a percentage of income; time-invariant, entered as a z-score (Percent of income (z-scored)) & DOE LEAD 2022 & 3,144 counties & Energy-burden mechanism \\
Farm earnings share & Farm earnings as a share of total county earnings over the baseline period; time-invariant, entered as a z-score (Share (z-scored)) & BEA CAINC5N / USDA ERS & 3,116 counties & Agricultural-channel tests (Appendix A) \\
\addlinespace
\multicolumn{5}{l}{\textit{Auxiliary}} \\
County population & Census population; used for population weights in rating-area aggregation and per-capita denominators (Persons) & Census & 1996--2023; 87,919 obs. & Weighting and aggregation \\
Net migration rate & Net IRS filer migration as a share of non-migrant filers (Rate) & IRS SOI & 2012--2020; 27,899 obs. & Selection bound (Appendix A) \\
Dew-point temperature & Annual mean dew-point temperature from 4km PRISM grids aggregated to counties (Degrees F) & PRISM & 2009--2023; 46,620 obs. & Humidity robustness \\
Uninsured rate & Share of the under-65 population without health insurance (Share) & Census SAHIE & 2012--2022; 34,543 obs. & Debt measurement critique \\
\end{longtable}
\vspace{-0.6em}\begin{minipage}{\textwidth}\footnotesize
Coverage columns are computed from the delivered panels: `Years' spans non-missing observations, `Obs.' counts them, and `Counties' is the number of distinct counties (states, for the state panel) with at least one non-missing value. Time-invariant moderators report a single county count and no year range. County dollar series are in 2023 USD; state dollar series are in 2025 USD. Shock indicators use a fixed national 80th-percentile threshold from the 1990--2000 baseline, so a county's shock status is comparable over time. The machine-readable CSV carries an additional `Column' field giving the exact variable name in the delivered panel.
\end{minipage}
\endgroup


Table~\ref{tab:e1t1} reports summary statistics for the estimation panel. The shock bins convey the treatment intensity directly: an extreme-heat year occurs in 24 percent of county-years and an extreme-cold year in 17 percent, while extreme drought - a Palmer reading at or below -4 - is rare, marking 2.3 percent of county-years unweighted and 4.5 percent when weighted by population. The outcome ledgers differ as much in level as in source. Mean county per-capita income is \$53,145 in 2023 dollars, with the middle eighty percent of counties between roughly \$39,000 and \$70,000; about 19 percent of adults in the average county have medical debt in collections, with wide dispersion across counties (5 percent at the tenth percentile, 32 percent at the ninetieth); the mean benchmark silver premium is \$374 per month; and Medicare spending averages \$10,359 per beneficiary, with 629 emergency-department visits per 1,000. County size is heavily right-skewed - mean population is about 103,000 but the median county has 26,000 residents, and median county employment is 10,800 against a mean of 48,000 - so the unweighted county regressions that follow describe the typical county, and population-weighted variants appear as robustness.

\begingroup
\footnotesize
\setlength{\tabcolsep}{4pt}
\sloppy
\begin{longtable}{>{\raggedright\arraybackslash}p{9.40cm} r >{\raggedright\arraybackslash}p{2.6cm} r >{\raggedright\arraybackslash}p{2.7cm} r}
\caption{Descriptive statistics, county panel 2011--2023}
\label{tab:e1t1}\\
\toprule
Variable & N & Mean (SD) & Median & P10--P90 & Missing \\
\midrule
\endfirsthead
\multicolumn{6}{l}{\textit{Table \thetable{} (continued)}} \\
\toprule
Variable & N & Mean (SD) & Median & P10--P90 & Missing \\
\midrule
\endhead
\midrule \multicolumn{6}{r}{\textit{continued on the next page}} \\
\endfoot
\bottomrule
\endlastfoot
\multicolumn{6}{l}{\textit{Climate exposure}} \\
Temperature, standard deviations from the county's 1990--2000 mean & 40,781 & 0.886 (1.07) & 1.02 & -0.743--2.15 & 2.58\% \\
Precipitation, standard deviations from the county's 1990--2000 mean & 40,781 & 0.178 (1.36) & 0.0977 & -1.39--1.90 & 2.58\% \\
Extreme heat (cooling degree days above the national 80th percentile) & 40,781 & 0.243 (0.429) & 0.00 & 0.00--1.00 & 2.58\% \\
Extreme cold (heating degree days above the national 80th percentile) & 40,781 & 0.174 (0.379) & 0.00 & 0.00--1.00 & 2.58\% \\
Palmer Drought Severity Index, annual mean & 40,391 & 0.320 (2.17) & 0.196 & -2.51--3.15 & 3.52\% \\
Extreme drought (Palmer index at or below -4) & 40,781 & 0.0230 (0.150) & 0.00 & 0.00--0.00 & 2.58\% \\
\addlinespace
\multicolumn{6}{l}{\textit{Health and financial outcomes}} \\
Share of adults with medical debt in collections & 38,810 & 0.188 (0.0996) & 0.185 & 0.0540--0.322 & 7.29\% \\
Median medical debt in collections (2023 USD) & 26,618 & 955 (389) & 892 & 521--1,464 & 36.4\% \\
Benchmark silver premium, monthly (2023 USD) & 30,690 & 374 (119) & 346 & 247--534 & 26.7\% \\
Lowest bronze premium, monthly (2023 USD) & 30,690 & 300 (87.7) & 283 & 204--411 & 26.7\% \\
Hospital bad debt, annual total (2023 USD) & 31,437 & 6,808,145 (21,670,523) & 1,710,004 & 272,891--13,877,415 & 24.9\% \\
Hospital charity care, annual total (2023 USD) & 31,437 & 10,059,331 (48,847,670) & 936,409 & 0.00--20,387,260 & 24.9\% \\
Uninsured share of the under-65 population & 34,543 & 0.118 (0.0577) & 0.110 & 0.0530--0.193 & 17.5\% \\
\addlinespace
\multicolumn{6}{l}{\textit{Household and local economy}} \\
Per-capita personal income (2023 USD) & 40,937 & 53,145 (16,735) & 50,565 & 39,087--70,051 & 2.21\% \\
Median household income (2023 USD) & 40,148 & 62,433 (16,327) & 60,255 & 44,721--82,113 & 4.10\% \\
Civilian employment (persons) & 40,154 & 48,068 (157,164) & 10,773 & 2,092--98,364 & 4.08\% \\
Household income, credit-bureau file (2023 USD) & 37,681 & 81,313 (20,231) & 77,720 & 61,037--104,661 & 9.99\% \\
County population & 40,840 & 103,466 (329,584) & 25,776 & 5,019--210,582 & 2.44\% \\
\end{longtable}
\vspace{-0.6em}\begin{minipage}{\textwidth}\footnotesize
One row per variable, pooled over county-years in the estimation panel. N counts county-years with a value and `Missing' the share of panel rows without one; `P10--P90' brackets the middle eighty percent of counties. The two shock indicators take the value 1 in a shock year and 0 otherwise, so their means are the share of county-years flagged. Means and standard deviations are given to three significant figures and dollar series are in 2023 dollars. County size is heavily right-skewed -- the mean county has about four times the employment of the median -- so these unweighted means describe the typical county rather than the typical resident.
\end{minipage}
\endgroup


Two county-level moderators enter this essay's mechanism analysis: the climate-exposed non-farm labor share and the DOE LEAD energy-burden z-score. The Social Vulnerability Index and the hospital-side moderators are the subject of the third essay and appear here for reference.

Table~\ref{tab:e1t0b} defines all variables the essay uses, grouped into shocks, outcomes, moderators, and auxiliary series. The shock measures are binary and defined against a fixed national threshold - the 80th percentile of the 1990--2000 baseline distribution - rather than a within-county quantile, so a county's shock status is comparable both over time and across counties, and counties with mild climates can go the whole panel without a shock. The two labor moderators are time-invariant and enter as z-scores, so the interaction coefficients reported later are per-standard-deviation gradients within a shock year rather than levels of employment loss. Employment enters in logs because the levels specification weights large counties heavily.

All fixed-effects models include county and year effects (state and year in the state panel), with standard errors clustered at the state level. Following \citet{abadie2023clustering}, the clustering level should match the level at which treatment is effectively assigned and correlated, and climate shocks are spatially correlated well beyond county borders. The choice matters for inference but not for the point estimates: clustering at the county level would be extremely non-conservative on headline p-values, while Conley standard errors with a 200-km kernel essentially reproduce state clustering and leave the headline estimates intact (Appendix A). \citep{conley1999gmm}

\FloatBarrier
\section{Medicare morbidity and utilization}

This essay's question concerns healthcare costs and therefore I looked for data sources where these costs were recorded in as much detail as possible at the population level. The CMS Geographic Variation files cover the Medicare population (those 65 and over and the disabled). These are administratively complete and their standardized payments strip out price variation, so utilization responses can be observed without churn in insurance composition. In this section, I study the morbidity channel in Medicare claims directly and do not extend it to working-age households.

I estimate distributed-lag fixed-effects models with county and year effects and state-clustered standard errors, using within-county variation in the heat, cold, and air-quality bins at lags of zero to two years. The sample is 30,641 county-years across 3,124 counties over 2014-2023, and the outcomes are standardized per-beneficiary spending and emergency-department visits per 1,000 beneficiaries (Table~\ref{tab:e1t6}).

\begingroup
\footnotesize
\setlength{\tabcolsep}{4pt}
\sloppy
\begin{longtable}{>{\raggedright\arraybackslash}p{12.70cm} >{\raggedright\arraybackslash}p{2cm} l r r r}
\caption{Medicare spending and utilization responses to climate and air-quality shocks}
\label{tab:e1t6}\\
\toprule
Hazard & Timing & Estimate & Std. error & p & Percent of mean \\
\midrule
\endfirsthead
\multicolumn{6}{l}{\textit{Table \thetable{} (continued)}} \\
\toprule
Hazard & Timing & Estimate & Std. error & p & Percent of mean \\
\midrule
\endhead
\midrule \multicolumn{6}{r}{\textit{continued on the next page}} \\
\endfoot
\bottomrule
\endlastfoot
\multicolumn{6}{l}{\textit{Emergency department visits per 1,000}} \\
Extreme heat & Same year & 7.77 & 2.71 & 0.006 & 1.24\% \\
Extreme heat & 1 year & 9.44 & 2.38 & \textless{}0.001 & 1.50\% \\
Extreme heat & 2 years & 1.99 & 2.52 & 0.433 & 0.317\% \\
\addlinespace
\multicolumn{6}{l}{\textit{Standardized spending per beneficiary}} \\
Extreme heat & Same year & \$112 & 43.5 & 0.013 & 1.08\% \\
Extreme heat & 1 year & \$176 & 52.2 & 0.002 & 1.69\% \\
Extreme heat & 2 years & \$75.3 & 24.3 & 0.003 & 0.727\% \\
\addlinespace
\multicolumn{6}{l}{\textit{Emergency department visits per 1,000}} \\
Extreme cold & Same year & 1.06 & 3.04 & 0.728 & 0.169\% \\
Extreme cold & 1 year & 2.46 & 2.07 & 0.241 & 0.391\% \\
Extreme cold & 2 years & 9.01 & 2.80 & 0.002 & 1.43\% \\
\addlinespace
\multicolumn{6}{l}{\textit{Standardized spending per beneficiary}} \\
Extreme cold & Same year & \$38.2 & 34.1 & 0.268 & 0.369\% \\
Extreme cold & 1 year & \$39.2 & 30.0 & 0.197 & 0.378\% \\
Extreme cold & 2 years & \$86.8 & 31.5 & 0.008 & 0.838\% \\
\addlinespace
\multicolumn{6}{l}{\textit{Emergency department visits per 1,000}} \\
Poor air quality & Same year & 5.00 & 1.25 & \textless{}0.001 & 0.795\% \\
Poor air quality & 1 year & 3.64 & 1.07 & 0.001 & 0.578\% \\
Poor air quality & 2 years & 2.83 & 1.21 & 0.023 & 0.450\% \\
\addlinespace
\multicolumn{6}{l}{\textit{Standardized spending per beneficiary}} \\
Poor air quality & Same year & \$19.0 & 17.8 & 0.290 & 0.184\% \\
Poor air quality & 1 year & \$2.10 & 20.7 & 0.919 & 0.0203\% \\
Poor air quality & 2 years & -\$14.1 & 17.8 & 0.434 & -0.136\% \\
\addlinespace
\multicolumn{6}{l}{\textit{Emergency department visits per 1,000}} \\
Extreme drought & Same year & 7.73 & 4.27 & 0.076 & 1.23\% \\
Extreme drought & 1 year & -4.19 & 4.93 & 0.399 & -0.667\% \\
Extreme drought & 2 years & -4.23 & 5.84 & 0.472 & -0.672\% \\
\addlinespace
\multicolumn{6}{l}{\textit{Standardized spending per beneficiary}} \\
Extreme drought & Same year & \$118 & 72.2 & 0.108 & 1.14\% \\
Extreme drought & 1 year & \$21.7 & 53.3 & 0.686 & 0.210\% \\
Extreme drought & 2 years & -\$9.61 & 47.3 & 0.840 & -0.0928\% \\
\end{longtable}
\vspace{-0.6em}\begin{minipage}{\textwidth}\footnotesize
County-year panel of CMS Geographic Variation data, 2014--2023, for the 65-plus and disabled Medicare population; county and year fixed effects, standard errors clustered by state. `Percent of mean' expresses each estimate against its own baseline: 10,359 dollars of standardized spending per beneficiary and 629 emergency department visits per 1,000 beneficiaries.
\end{minipage}
\endgroup


\begin{equation}
Y_{it} \;=\; \alpha_i \;+\; \gamma_t \;+\; \sum_{\ell=0}^{2} \beta_\ell \, S_{i,t-\ell}
\;+\; X_{it}'\delta \;+\; \varepsilon_{it}
\label{eq:dlag}
\end{equation}

Extreme heat raises standardized Medicare spending by \$112 per beneficiary in the shock year (p = 0.013), \$176 the following year (p = 0.002), and \$75 two years out (p = 0.003). From a baseline mean of \$10,359 per beneficiary, these represent increases of 1.1, 1.7, and 0.7 percent respectively. Emergency-department visits rise by 7.8 per 1,000 beneficiaries in the shock year (p = 0.006) and 9.4 the following year (p = 0.0002); from a baseline mean of 629 per 1,000, representing increases of 1.2 and 1.5 percent. The lag structure is informative as the largest response arrives the year after exposure. This confirms that the costs of extreme weather are rarely instantaneous and reveal themselves over time.

Cold, on the other hand, operates more slowly: spending rises by \$87 per beneficiary and emergency visits by 9.0 per 1,000 at the two-year lag. Air-quality shocks raise emergency visits by 5.0 per 1,000 contemporaneously (p = 0.0002), 3.6 at one year, and 2.8 at two, in the 1,194 counties with an operating monitor, which are disproportionately urban. Figure~\ref{fig:e1f4} plots the dynamic profiles by hazard.

\begin{figure}[htbp]
\centering
\caption{Medicare spending and emergency-department responses by hazard and years since the shock.}
\label{fig:e1f4}
\includegraphics[width=0.95\linewidth]{../../../Analysis/mechanism/plots/fig_medicare_morbidity.png}
\end{figure}

These estimates need to be checked for robustness in five ways which I detail here. The first is outcome selection: with several utilization measures, one could appear significant by chance. Following \citet{anderson2008multiple}, I collapse the utilization outcomes into a single standardized summary index, and the index is significant at the same time lags where the component outcomes are (heat at one lag, p = 0.007; cold at two lags, p = 0.001). The second is multiple testing across the many hazard-outcome-lag cells. Under a sharpened Benjamini-Krieger-Yekutieli false-discovery correction, the heat response at one year and the contemporaneous air-quality response survive at q < 0.05. \citep{benjamini2006adaptive} The third is spatial correlation:  climate shocks do not stop at state borders and neighboring counties may be exposed to correlated shocks. The Conley standard errors, which allow correlation within a distance kernel, are tighter than the state-clustered ones. The fourth is recurring treatment: counties enter and leave shock status repeatedly, which can contaminate the control comparisons in two-way fixed-effects lag models. For robustness, I use the estimator built for this kind of treatment path \citep{dechaisemartin2026intertemporal}, which confirms a heat spending effect of approximately \$80 at a two-year horizon. The fifth is the identifying variation: absorbing state-by-year shocks leaves the heat-spending response at one lag imprecise (+\$34, p = 0.15), so that estimate rests on variation across states within a year as well as within them, though the emergency-visit responses are unaffected. The last is attribution, as the channel through which these shocks flow needs checking. The estimates reproduce \citet{deryugina2019mortality} within this panel.

Turning to the labor market, the lagged burden of climate shocks operates substantially outside the agricultural mechanism. I interact the heat-shock indicator (a county-year in which cooling degree days reach the national 80th percentile of the 1990-2000 baseline distribution) with two time-invariant county characteristics that capture exposure outside the farm sector: the baseline share of employment in climate-exposed non-farm industries (construction, manufacturing, transportation, and utilities; mean 24.9 percent, standard deviation 6.1 percentage points) and household energy burden (mean 3.4 percent of income, standard deviation 1.1 percentage points). Both gradients are negative. In a heat-shock year, a county one standard deviation above the mean in exposed non-farm employment share loses an additional 0.52 percent of employment (p = 0.006), and a county one standard deviation above the mean in energy burden loses an additional 0.84 percent (p = 0.005); at the median county's 10,773 workers, these amount to roughly 56 and 90 jobs.

Each gradient holds under the specification that most directly threatens it. Division-by-year fixed effects, which absorb regional shocks common to counties within a census division, leave both essentially unchanged (-0.0042, p = 0.015; -0.0078, p = 0.002). The labor-share gradient holds under Conley standard errors allowing for spatial correlation within 200 kilometers (p = 0.033), and the energy-burden gradient holds in a joint horse-race against agricultural dependence, the exposed-industry share, social vulnerability, and baseline climate (-0.0065, p = 0.020). These interactions describe how the employment response to heat varies with county exposure; they are not estimates of its average level. The analogous labor income margin is not supported by our data.

A note before going deeper into the financial impacts: the Medicare population is not the population of working-age individuals represented in credit bureau debt or of the household economic outcomes in Section 6. I am aiming to triangulate by directly measuring morbidity and household costs after climate shocks where possible, which in no way claims an identified causal chain.

\FloatBarrier
\section{Financial-ledger responses: debt and ACA premiums}

Which institutions record or price these climate-related costs? I examine two sources with different functions: credit-bureau medical debt, which records unpaid bills, and ACA benchmark premiums, which price expected claims. To summarize, the harm is recorded partially and with lags in debt, and it is not priced locally in any coherent way in premiums (Table~\ref{tab:e1t7}; Figure~\ref{fig:e1f5}).

\begin{table}[htbp]
\centering
\footnotesize
\setlength{\tabcolsep}{4pt}
\sloppy
\caption{Institutional ledgers compared: response, baseline, and response as a share of baseline}
\label{tab:e1t7}
\begin{tabularx}{\textwidth}{>{\raggedright\arraybackslash}p{3cm} >{\raggedright\arraybackslash}X r r r r}
\toprule
Ledger & Response & Estimate & p & Baseline & Percent of baseline \\
\midrule
Medicare (administrative) & Standardized spending per beneficiary, heat at 1-year lag & \$176 & 0.002 & \$10,359 & 1.69\% \\
Medicare (administrative) & Emergency visits per 1,000, heat at 1-year lag & 9.44 & \textless{}0.001 & 629 & 1.50\% \\
Credit bureau (state) & Medical debt share, cold at 1-year lag & 1.35 pp & 0.012 & 18.8\% & 7.19\% \\
Credit bureau (county) & Medical debt share, cold at 1-year lag & -0.272 pp & 0.463 & 18.8\% & -1.44\% \\
Credit bureau (county) & Medical debt share, drought at 2-year lag & 0.581 pp & 0.024 & 18.8\% & 3.09\% \\
ACA premiums (rating area) & Benchmark silver premium, drought at 2-year lag & \$3.13/mo & 0.236 & \$366/mo & 0.855\% \\
\bottomrule
\end{tabularx}
\begin{minipage}{\textwidth}\vspace{0.5em}\footnotesize
Each ledger records a different population under different rules, so responses are expressed as a percentage of that ledger's own mean to make them comparable. Medicare estimates cover 2014--2023 for the 65-plus and disabled population; the ACA cell is estimated on the rating-area pass-through panel. The two county debt cells are estimated as the county mirror of the state specification, because the county pipeline carries no binary cold-z shock and reports drought through a continuous index -- no committed county output contains them. Their values differ from figures quoted in earlier drafts and supersede them. All specifications use county (or state) and year fixed effects with state-clustered standard errors.
\end{minipage}
\end{table}


\begin{figure}[htbp]
\centering
\caption{Climate-shock responses across institutional ledgers, each expressed as a percent of the mean of that ledger.}
\label{fig:e1f5}
\includegraphics[width=0.95\linewidth]{../../../Analysis/mediation/fig_institutional_ledgers.png}
\end{figure}

Cold shocks raise the medical-debt share by 0.85 to 1.35 percentage points at a one-year lag in the state panel (p = 0.012 to 0.044 across baseline choices); the county mirror of the same specification does not reproduce it (-0.27 points, p = 0.46; Table~\ref{tab:e1t7}). From a mean of about 19 percent of adults with medical debt in collections, 1.35 points is a relative increase of roughly 7 percent. These are lagged associations from distributed-lag fixed-effects models rather than causal event estimates explored in Appendix A. The response survives a stricter 90th-percentile shock definition (1.38 points, p = 0.003) and humidity adjustment. The divergence between the two panels is itself informative about the ledger: bureau debt is recorded only after insurance, billing, and credit-file filters, and those filters differ enough across counties within a state that the association visible in state aggregates does not survive disaggregation.

Drought's debt response is weaker and depends on the level of aggregation: the debt share rises by 0.58 percentage points at a two-year lag in the county panel (p = 0.02), an estimate that does not survive the sample restriction those controls impose - on the identical sample it falls to 0.09 points (p = 0.76) before any control is added, while the state-level specification gives an insignificant 0.72 points (p = 0.18). Therefore the evidence remains mixed.

Bureau debt behaves this way because three filters stand between exposure and recording in data sources.  An insurance relationship has to exist, there has to be a billed encounter, and a credit file has to exist on an individual. Incidence failing these filters never enters the data source, and reporting-rule changes over the window alter what is recorded. The third essay develops the full measurement exploration; for present purposes it suffices to note that debt is a selective data source.

The premium test is specified by the institutions that set premiums. Rates are set at the rating-area level and reviewed by states; plan-year-t rates are filed around the middle of year t-1 on claims experience through roughly t-2, then locked for the plan year. Only lagged shocks can therefore be in the insurer's information set, and I take the two-year lag as primary. I estimate at the two levels the institutions themselves use -- rating-area-by-year as the primary specification (rating-area and state-by-year fixed effects, population-weighted, state-clustered) and state-by-year as the secondary -- and display county-level specifications only as a sanity check. This is because roughly 86 percent of premium variance is state-by-year, so a county-plus-year-fixed-effects premium regression is confounded.

\begin{equation}
P_{rt} \;=\; \alpha_r \;+\; \gamma_{s(r),t} \;+\; \theta \, \bar{S}_{r,t-2} \;+\; \varepsilon_{rt}
\label{eq:passthrough}
\end{equation}

The results show sign instability across the levels at which prices are actually set. The cold coefficient at the two-year lag runs +\$12.6 at the rating-area level and -\$16.7 at the state level; heat runs -\$10.5 and +\$93; drought is null at both. County-level estimates are reported as a sanity check only, since roughly 86 percent of premium variance is state-by-year. A genuine price response would keep its sign as aggregation moves through the levels at which prices are actually set, so this pattern indicates no coherent local pass-through. The one large coefficient - the between-state heat estimate of \$54 to \$93 per month - is approximately ten times too large for a claims channel, and the state-level cold coefficient has the wrong sign relative to this essay's own Medicare result, in which cold raises spending.

Finding no effect can mean two very different things. It can mean premiums genuinely do not respond to climate shocks, or it can mean the response is real but too small for this panel to detect. A coefficient near zero cannot tell them apart on its own. To separate them, I ask a more demanding question: how large a premium response would the data have caught if one existed? Anything larger than that threshold can be ruled out, and anything smaller remains possible.

The natural yardstick is the premium increase that insurers would have to charge if they fully priced the health costs the earlier sections measure. Working from the Medicare morbidity estimates, that figure is \$9.30 to \$14.63 per member per month. This is the benchmark for full pass-through: if climate shocks raise medical costs by this much and insurers pass all of it on, premiums should move by roughly this amount. The shock variables are measured as the share of a rating area's population exposed in a given year, so each coefficient is the premium change for an area that goes from none of its population exposed to all of it. That puts the estimates and the benchmark in the same units and makes the comparison meaningful.

Drought gives a clear answer. The estimated response is \$3.13 per member-month, with a standard error of \$2.60 - small, and not statistically distinguishable from zero. More importantly, the estimate is precise enough to rule things out: the data reject any response above \$7.41 per member-month. That threshold sits below the full pass-through benchmark, at 51 to 80 percent of it depending on which end of the benchmark range one takes. Insurers therefore cannot be passing drought-related health costs into premiums at anything close to the scale those costs imply. This is a real finding about how the market prices risk.

Heat and cold give no answer at all. Both estimates are far less precise than the drought estimate. The smallest response the data can rule out is \$24.61 per member-month for heat and \$22.04 for cold. Both thresholds are larger than the full pass-through benchmark itself, which means the test cannot rule out a response as large as full pass-through - but it equally cannot rule out a response of zero. For these two hazards the exercise is uninformative. It is not evidence that heat and cold costs are passed through, and it should not be read that way; it is a statement about how much we can observe in this panel. What the data do pin down for heat and cold is an upper limit: whatever the response is, it is no larger than about 6 to 7 percent of the population-weighted mean benchmark premium in the rating-area estimation sample, \$366 per month.

For drought alone the data rule out a response above \$7.41 per member per month; for heat and cold the bounds are too wide to be informative in either direction.

The premium null has a direct implication for the debt results above. If premiums do not move when climate shocks - specifically drought - hit, then the increase in medical debt cannot be running through premiums, because there is no price change for it to run through. The data bear this out. When I re-estimate each shock--debt relationship with premiums included as a control, nearly all of the original effect remains: 98.7 percent for drought at a two-year lag, and 92.2 percent for cold at a one-year lag. Remaining shock-by-lag estimates for other exposures are statistically non-significant and uninformative. Controlling for premiums leaves the debt results essentially unchanged.

The null result can be explained institutionally. A single statewide risk pool, a geographic rating factor limited to unit costs, and federal risk adjustment under Section 153 each attenuate any local claims signal before it reaches a local premium. \citep{kautter2014hhshcc} In this sense, regulated premium-setting does not locally price the environmental lag structure documented in this essay.

A final contrast closes the comparison between data sources. State per-capita health spending shows one climate-shock association, with extreme cold (+\$205, p = 0.014), and tracks unemployment (p = 0.06) but not income (p = 0.40). This is why the essay looks for climate costs in lagged household and administrative data sources: while the costs are real (Sections 4 and 6 measure them),  the contemporaneous system-level spending aggregates do not reveal them.

\FloatBarrier
\section{Household economic capacity: income and employment}

The premium and debt results presuppose a financing margin: unpaid medical bills accumulate where household budgets cannot absorb them, so the recording patterns of Section 5 are most interpretable alongside direct evidence that climate shocks strain local incomes. The drought coefficient on county per-capita income is stable between -\$99 and -\$132 per unit of drought severity across estimation windows beginning in 1990, 2000, and 2011, with precision improving as the window lengthens; in the full window the contemporaneous term is also significant (-\$149, p = <0.01).

The 2012 drought provides an event-based complement (Appendix A). The raw contrast between the 139 first-onset counties and 2,534 never-exposed controls is -\$1,311 - about 2.5 percent of mean income, averaged over the onset year and the eleven years that follow, and it survives wild-cluster bootstrap (p = 0.036) and randomization inference (p = 0.0075). Decomposing county income we find that treated-county farm earnings spiked to \$4,339 per resident in 2011 - the commodity-price peak - and the farm component of the contrast falls from -\$907 to -\$14 once the baseline pools 2009--2011, consistent with mean reversion rather than drought damage. The component stable under every baseline choice is a nonfarm decline of \$261 to \$414, roughly 0.5-0.8 of a percent of mean income. The decomposition shows that in agricultural counties the income sources embed the commodity-price cycle, and event contrasts need to separate farm and nonfarm income to get a fuller picture.

Employment shows the same event-specific pattern with less stability: the 2012 event reduced employment by roughly 2,000 jobs per treated county - about 4 percent of mean county employment (Appendix A.3). That estimate is materially more fragile than the income one: it attenuates by about 58 percent under the doubly-robust estimator (-871, SE 433) and reverses sign in the pooled estimator (+2,609, SE 2,245). Because the employment series begins in 2011, no pre-onset leads exist either, so I report it as event-specific and secondary. Taken together, the capacity results indicate that climate shocks strain the budgets from which medical bills are paid; they support the plausibility of the debt responses in Section 5 without claiming an identified chain from income to debt.

\FloatBarrier
\section{Interpretation and limitations}

Taken together, the results say the following: after climate shocks, morbidity is directly measurable in administrative health data; the regulated premium pricing does not coherently price it locally; and household economic capacity strains. The connective claim is triangulation across ledgers that observe different populations under different rules. The essay does not identify a propagation chain from one source to another.

The external-validity claims we can take out of this essay are very specific. The Medicare evidence covers the 65-plus and disabled population over 2014--2023. The premium results cover the ACA individual market in the marketplace era, and the tight equivalence bound applies to drought alone. The 2012 event contrast is an intention-to-treat quantity for its treated geography; the pooled estimator indicates it does not generalize, and the decomposition of Appendix A bounds its causal component.

Three measurement caveats qualify the results. Credit-bureau debt records hardship only after insurance, billing, and credit-file filters, so it understates hardship where coverage and credit access are thin - this is tackled in the third essay. Premiums are a regulated price, so a premium null is informative about institutions rather than evidence that costs are absent. And the ACS-based moderators are smoothed over time, limiting the ability to detect fast demographic adjustment.

Two future extensions would sharpen our understanding of climate exposure and financial dynamics. Working-age administrative utilization data - all-payer claims - would test the morbidity channel in the same population as the credit bureau debt. Patient-flow measures of hospital exposure would replace the location-county assignment used on the provider side.

\FloatBarrier
\section{Conclusion}

This essay asked what healthcare costs follow climate shocks, and which financial institutions record or price them. Three answers emerge: Medicare records show extreme heat raising standardized spending by \$112 per beneficiary in the shock year and \$176 the next, with emergency visits up 8-9 per 1,000, among the 65-plus and disabled population. ACA premiums show no coherent local pass-through of any of the climate shocks. Lastly, these unpriced costs are absorbed in strained household budgets, as drought lowers local incomes.

These insights inform the second and third essays, which examine how these costs evolve - whether they persist or compound - and how they are distributed. We can only detect a full picture of climate-driven financial costs where administration is universal. Medicare and the BEA income accounts observe different populations under different rules, so no single ledger sees the whole burden. For climate adaptation and health-finance policy, the immediate implication is about measurement: no single source observes the entire cost burden, and therefore estimation and modeling choices should take this into account.

\FloatBarrier
\FloatBarrier
\appendix
\section{The 2012 drought natural experiment}

\subsection{The 2012 drought design}

Identification needs an exposure that is sharp, unanticipated, and spatially delimited, and the 2012 drought onset provides all three. I restrict treatment to counties experiencing their first panel-era extreme drought in 2012, and I restrict the comparison group to counties never exposed to extreme drought in the window, which avoids the problematic comparisons that arise in staggered designs when already-treated units serve as controls. The restriction yields 139 treated counties - in Georgia, the Mountain West, and the Plains - and 2,534 never-exposed controls; Figure~\ref{fig:e1f1} maps the two groups and Table~\ref{tab:e1t2} reports balance. Treated counties are smaller and more rural than the average control, which motivates both the covariate-adjusted doubly-robust estimator and the two-decade pre-trend test reported below.

\begin{table}[htbp]
\centering
\footnotesize
\setlength{\tabcolsep}{4pt}
\sloppy
\caption{Pre-treatment (2011) balance: 2012 first-onset counties versus never-exposed counties}
\label{tab:e1t2}
\begin{tabularx}{\textwidth}{>{\raggedright\arraybackslash}X r r r r}
\toprule
Variable & Treated (139) & Never-exposed (2,534) & Difference & Norm. diff. \\
\midrule
Per-capita income (2023 USD) & 48,488 & 49,297 & -809 & -0.064 \\
Civilian employment & 16,641 & 44,456 & -27,815 & -0.314 \\
Population & 37,812 & 96,741 & -58,929 & -0.309 \\
Medical debt share & 0.231 & 0.212 & 0.020 & 0.199 \\
Median household income (2023 USD) & 74,362 & 78,751 & -4,389 & -0.250 \\
\bottomrule
\end{tabularx}
\begin{minipage}{\textwidth}\vspace{0.5em}\footnotesize
Treated counties are the 139 whose first extreme-drought year (PDSI at or below -4) in the 2011--2023 panel is 2012; never-exposed counties are the 2,534 with no extreme-drought year in the window. Means are taken in 2011, the single pre-period of the two-by-two design. The normalized difference is the treated-control gap divided by the square root of the average of the two group variances; values above 0.25 in absolute terms are conventionally treated as substantial imbalance, and motivate the doubly-robust estimator reported alongside the raw contrast.
\end{minipage}
\end{table}


\begin{figure}[htbp]
\centering
\caption{Counties with their first extreme-drought onset in 2012 (treated) against counties never exposed to extreme drought over the panel.}
\label{fig:e1f1}
\includegraphics[width=0.95\linewidth]{../../../Analysis/did/fig_treated_map.png}
\end{figure}

I estimate a two-by-two difference-in-differences of each outcome on the interaction of treatment and a post-2012 indicator, with county and year fixed effects and state-clustered standard errors. The estimand is the intention-to-treat effect of first extreme-drought onset for the 139 treated counties over 2011--2023, relative to never-exposed counties, identified under the assumption that the two groups would have trended in parallel absent the onset. Three threats to that assumption structure the falsification agenda reported in A.3 and summarized in Table~\ref{tab:e1t4}: selection into drought-prone geography, which would appear as differential pre-trends; dependence on a single state or region, addressed with leave-one-state-out estimates and placebo onsets; and drought spillovers into neighboring control counties, which would violate the no-interference assumption. Because the treated counties fall in only 17 states, small-cluster distortions are a first-order concern, and I report analytic, wild-cluster-bootstrap, and randomization-inference p-values side by side for every headline estimate.

\begin{equation}
Y_{it} \;=\; \alpha_i \;+\; \gamma_t \;+\; \tau \,\bigl(D_i \times \mathrm{Post}_t\bigr)
\;+\; \varepsilon_{it}
\label{eq:did2x2}
\end{equation}

\FloatBarrier
\subsection{The event contrast and its decomposition}

I begin with the raw event contrast. Comparing the 139 counties whose first extreme-drought onset came in 2012 with the 2,534 counties never exposed over the panel, in a two-by-two difference-in-differences design with county and year fixed effects, per-capita income in treated counties averages \$1,311 below never-exposed controls over the onset year and the eleven years that follow (analytic p = 0.028). From a mean county per-capita income of \$53,145 (2023 dollars; Table~\ref{tab:e1t1}), this is about 2.5 percent. The contrast is an intention-to-treat quantity for a specific event and geography - Georgia, the Mountain West, and the Plains. Two immediate objections do not overturn it, and Table~\ref{tab:e1t4} reports both alongside the rest of the falsification suite. Because the treated counties fall in only 17 states, conventional cluster-robust inference can overstate precision; a wild-cluster bootstrap with Webb weights yields p = 0.036 with a confidence interval of [-\$2,911, -\$139], and randomization inference yields p = 0.0075. \citep{cameron2008bootstrap,mackinnonwebb2018wildbootstrap,young2019channelingfisher} Because treated counties are more rural and more agricultural than the average control, the comparison might load observable composition; a doubly-robust estimator conditioning on county characteristics returns -\$1,451 (SE 515). One feature of the design survives both checks unchanged: like the two-by-two itself, the doubly-robust estimator measures every post-period gap against the single pre-year 2011. The rest of this section examines that baseline.

The within-panel event study around the 2012 onset (Figure~\ref{fig:e1f7}) shows the income gap opening immediately - \$1,561 in the onset year (p = 0.008) and \$1,609 the year after (p < 0.001) - narrowing through 2015, and re-widening to between \$1,300 and \$1,900 in the later post years; the two-by-two contrast of -\$1,311 is the average of these twelve post-year gaps. The employment gap, by contrast, builds with time rather than fading. Extending the income series backward, however, shows that the pre-onset years resemble the post-onset years: the 2008--2010 gaps relative to 2011 are -\$1,459 to -\$1,591, statistically indistinguishable from zero (p = 0.17 to 0.32) but of the same magnitude as the treatment-era gaps. The pooled contrast is correspondingly sensitive to the pre-period: widening the baseline from 2011 alone to 2009--2011 moves it from -\$1,311 (p = 0.028) to -\$285 (p = 0.64); Table~\ref{tab:e1t8} reports the full grid of contrasts across pre-periods and income components. Two readings are possible - treated counties may have been on a different trajectory throughout, or 2011 itself may have been an unusual year for treated counties. The county income accounts distinguish these directly.

\begin{table}[htbp]
\centering
\footnotesize
\setlength{\tabcolsep}{4pt}
\sloppy
\caption{Sensitivity of the 2012 drought income contrast to the pre-period baseline}
\label{tab:e1t8}
\begin{tabularx}{\textwidth}{>{\raggedright\arraybackslash}X r r r r}
\toprule
Pre-period baseline & Estimate & Std. error & p & County-years \\
\midrule
\multicolumn{5}{l}{\textit{Total per-capita income}} \\
2011 only & -\$1,311 & 577 & 0.028 & 34,086 \\
2010--2011 & -\$515 & 402 & 0.206 & 36,708 \\
2009--2011 & -\$286 & 599 & 0.636 & 39,330 \\
2007--2011 & -\$231 & 716 & 0.748 & 44,574 \\
2002--2011 & -\$419 & 1,059 & 0.694 & 57,684 \\
\addlinespace
\multicolumn{5}{l}{\textit{Farm earnings component}} \\
2011 only & -\$907 & 590 & 0.131 & 34,006 \\
2010--2011 & -\$246 & 216 & 0.260 & 36,618 \\
2009--2011 & -\$14.0 & 243 & 0.954 & 39,230 \\
2007--2011 & \$104 & 311 & 0.739 & 44,454 \\
2002--2011 & \$86.4 & 382 & 0.822 & 57,514 \\
\addlinespace
\multicolumn{5}{l}{\textit{Nonfarm income component}} \\
2011 only & -\$394 & 398 & 0.328 & 34,006 \\
2010--2011 & -\$261 & 437 & 0.552 & 36,618 \\
2009--2011 & -\$261 & 502 & 0.605 & 39,230 \\
2007--2011 & -\$286 & 521 & 0.586 & 44,454 \\
2002--2011 & -\$414 & 747 & 0.582 & 57,514 \\
\bottomrule
\end{tabularx}
\begin{minipage}{\textwidth}\vspace{0.5em}\footnotesize
Each row re-estimates the same two-by-two contrast -- 139 counties with their first extreme drought in 2012 against 2,534 never-exposed counties -- changing only which years form the pre-period against which post-2012 outcomes are measured. Reading down the first panel, the headline contrast shrinks from roughly \$1,300 to under \$300 as soon as the baseline extends beyond 2011 alone. The lower panels locate that instability: the farm component collapses from -\$907 to near zero, because 2011 was the peak of the commodity-price cycle, while the nonfarm component holds between -\$261 and -\$414 under every choice. County and year fixed effects with state-clustered standard errors throughout.
\end{minipage}
\end{table}


\begin{figure}[htbp]
\centering
\caption{Event-study panels around the 2012 drought onset: per-capita income, with twenty-one pre-treatment leads, and civilian employment.}
\label{fig:e1f7}
\includegraphics[width=0.95\linewidth]{../../../Analysis/plots/did/eventstudy_panels_Drought_2012.png}
\end{figure}

\begin{equation}
Y_{it} \;=\; \alpha_i \;+\; \gamma_t \;+\; \sum_{k \neq -1} \theta_k \,
\bigl(D_i \times \mathbf{1}\{t - 2012 = k\}\bigr) \;+\; \varepsilon_{it}
\label{eq:eventstudy}
\end{equation}

BEA county accounts separate farm earnings from other income, and the farm series identifies the source of the pattern (Figure~\ref{fig:e1f6}). Real farm earnings per treated-county resident averaged \$1,903 to \$2,438 over 2007--2010, then jumped to \$4,339 in 2011 - the year commodity prices peaked - while control-county farm income rose far less. Measured against that peak, the farm component accounts for roughly 85 percent of the 2008--2010 leads (-\$1,293 to -\$1,358 of the -\$1,459 to -\$1,591 totals). The farm component of the treatment contrast behaves accordingly: -\$907 (p = 0.13) against the 2011 baseline, and -\$14 (p = 0.95) once the baseline pools 2009--2011, consistent with mean reversion from the price peak rather than drought damage. The nonfarm component is steadier: -\$394, -\$261, -\$261, -\$286, and -\$414 across pre-periods beginning in 2011, 2010, 2009, 2007, and 2002 - sign-stable at 0.5 to 0.8 percent of mean income under every baseline choice, with onset-year gaps of -\$536 and -\$542 (p $\approx$ 0.08), though never significant at conventional levels. Taken together, the decomposition indicates that the raw contrast of -\$1,311 combines the reversion of farm income from its 2011 peak with a more modest, imprecisely estimated decline in nonfarm income. In agricultural counties the income ledger embeds the commodity-price cycle, and an event contrast inherits that cycle unless the two are separated.

\begin{figure}[htbp]
\centering
\caption{Farm versus nonfarm decomposition of the 2012 treated-control income gap, 2001--2023.}
\label{fig:e1f6}
\includegraphics[width=0.95\linewidth]{../../../Analysis/plots/did/decomposition_Drought_2012_farm_nonfarm.png}
\end{figure}

Section 6 of the main text reports the durable form of the drought--income relationship - the distributed-lag specification, which is stable across estimation windows and does not depend on any single baseline year. The contribution of this appendix is to establish what the event contrast measures: a commodity-price reversion in farm income combined with a modest nonfarm decline of 0.5 to 0.8 percent, the latter within the range of income effects in the climate-economy literature \citep{deschenes2007agricultural,deryugina2017fiscal}.

\FloatBarrier
\subsection{Falsification, external validity, and employment}

The identifying assumption is that treated and never-exposed counties would have trended in parallel absent the 2012 onset. Table~\ref{tab:e1t4} collects the tests of that assumption, each statistic taken from the committed robustness output rather than restated here, and three points deserve emphasis beyond what the table reports. First, the parallel-trends evidence is genuinely mixed: the differential linear pre-trend over the 21 years from 1990 to 2011 is -\$69 per year (SE 89, p = 0.44) and does not reject, but a coefficient of that size accumulates to roughly -\$1,450 over two decades - the magnitude of the raw event contrast - and an event-study joint Wald test does reject (F = 6.9, p < 0.001). A.2 locates that drift in the farm component of income, which is why this appendix reports the contrast under several baselines instead of resting on a single pre-trend statistic. Second, HonestDiD-style sensitivity analysis \citep{rambachan2023honestdid} cannot run on the 2012 cohort, which has no in-panel pre-period; run on the pooled cohorts, where a pre-period does exist, the income bounds include zero at every value of the smoothness parameter; the long BEA window and the farm/nonfarm decomposition substitute for it, and they also cover the falsification that future shocks should not predict past outcomes. Third, dropping Colorado or Nebraska lifts the analytic p-value to 0.075 and 0.057, which is the few-treated-cluster sensitivity that motivated the bootstrap rather than a separate finding. Three further checks are reported in full in the robustness suite: PRISM dew-point humidity leaves the headline coefficients essentially unchanged (the most sensitive candidate, the cold--debt coefficient, moves from 0.01363 to 0.01368), ACS controls for aging, tenure, and in-migration leave 94 to 104 percent of the debt and hospital effects intact, though the cold-income estimate retains only 58 percent, and stricter 90th-percentile shock definitions leave the conclusions unchanged.

\begin{table}[htbp]
\centering
\footnotesize
\setlength{\tabcolsep}{4pt}
\sloppy
\caption{Inference and falsification tests, 2012 drought event contrast (per-capita income)}
\label{tab:e1t4}
\begin{tabularx}{\textwidth}{>{\raggedright\arraybackslash}p{3.1cm} >{\raggedright\arraybackslash}X >{\raggedright\arraybackslash}p{4.2cm} >{\raggedright\arraybackslash}p{3.4cm}}
\toprule
Test & What it addresses & Statistic & Verdict \\
\midrule
Parallel pre-trends (BEA, 1990--2011) & Treated and control counties on different trajectories before 2012 & -\$69/yr (SE 89), p = 0.441; joint Wald F = 6.9, p \textless{}0.001 & Slope does not reject; joint test does. Drift is located in the farm component (A.2), not treated as settled by this test. \\
Wild-cluster bootstrap (Webb) & Only 17 treated states; analytic clustering overstates precision & p = 0.036; CI [-\$2,911, -\$139] & Survives; the bootstrap interval is the one cited in the text. \\
Randomization inference & Same, without asymptotic approximation & p = 0.007 & Survives. \\
Doubly-robust estimator & Treated counties differ in observable composition & -\$1,451 (SE 515); CI [-\$2,461, -\$441] & Survives; conditioning on observables does not erode the contrast. \\
Leave-one-treated-state-out & One influential treated state drives the result & ATT envelope [-\$1,687, -\$914] over 17 drops & Envelope stays inside the bootstrap interval; analytic significance flips on 2 drops (CO, NE). \\
Placebo onset years & The contrast would arise at any onset year & two-sided p = 0.009 (B = 1,000) & Estimate sits in the far tail of the placebo distribution. \\
Pre-window sensitivity & The estimate depends on the single 2011 pre-year & pre = 2011: -\$1,311 (p = 0.028); pre = 2009--2011: -\$286 (p = 0.636) & Not robust: the contrast is a function of the baseline year (see A.2). \\
\bottomrule
\end{tabularx}
\begin{minipage}{\textwidth}\vspace{0.5em}\footnotesize
All statistics are drawn from the robustness suite. The estimand throughout is the intention-to-treat effect of first extreme-drought onset in 2012 for the 139 treated counties, against 2,534 never-exposed counties, with county and year fixed effects and state-clustered standard errors. Further checks are reported in the text and appendices -- humidity, ACS demographic controls, threshold grids, and Conley spatial kernels. 
\end{minipage}
\end{table}


Drought in one county may depress income in its neighbors through shared labor and product markets, which would violate the no-interference assumption if control counties absorb part of the treatment. An own-versus-neighbor decomposition cannot separate the two here - own and adjacent-county exposure correlate at 0.94 to 0.97 - but two facts are identified. The neighbor-exposure block is jointly significant (p $\approx$ 0.006), and specifications adding neighbor exposure exceed the own-only baseline (income -157 versus -133 per unit of exposure; employment -1,090 versus -714). County coefficients therefore capture local exposure, and adjacent-county exposure adds a same-signed regional component that the local coefficient understates: spillovers amplify the estimate rather than confound it, and the headline is a lower bound on regional exposure. Spatial dependence in the errors is likewise handled: Conley standard errors with kernels out to 300 km leave the income result significant (worst case p = 0.008), and at 200 km the Conley and state-clustered p-values essentially coincide (0.0029 versus 0.0026).

Is the 2012 estimate a window onto droughts generally, or the fingerprint of one event? The panel contains later first-onset cohorts, so the question can be answered with estimates: I pool all first-onset cohorts in a staggered design and aggregate with the doubly-robust estimator of \citet{callaway2021did}, in the \citet{santanna2020drdid} specification. The pooled onset effect is -\$324 (SE 276), and the pooled long-run average is +\$350 (SE 585); both are nulls. The 2012 income effect is therefore event-specific even at onset, and I keep the event estimand and the pooled estimand separate throughout rather than averaging them. Several features of 2012 are consistent with the difference: it was an extreme event even among extremes; a first onset falls on counties with no recent adaptation experience, while later cohorts may be partially adapted; its timing within the agricultural calendar was unusually damaging; and the composition of comparison groups differs across cohorts. The data cannot adjudicate among these explanations, and I do not attempt to.

The 2012 event also reduced employment - by roughly 2,000 jobs per treated county (analytic p = 0.0001), about 4 percent of mean county employment of 48,068 - but this estimate is fragile in a way the income estimate is not. The fragility is not about inference: it clears the few-cluster bar (wild-cluster bootstrap p = 0.003; randomization inference p = 0.037), and its leave-one-state-out envelope of [-2,156, -1,854] never flips significance. Rather, it attenuates by about 58 percent under the doubly-robust estimator (-871, SE 433, with a confidence interval that barely excludes zero) and reverses sign in the pooled estimator (+2,609, SE 2,245) alongside positive pre-trends. I therefore report the employment decline as event-specific and secondary.

These tests defend a specific estimand: the intention-to-treat effect of the 2012 first onset for the treated geography, which the pooled cohorts confirm is event-specific by construction. The design does not identify the mechanism - agricultural losses, labor-demand contractions, and local-demand responses are all consistent with the reduced form - and the direct evidence on a health channel comes from the Medicare panel in Section 4. What this appendix establishes is what the event contrast measures; the durable form of the drought--income relationship is reported in Section 6 of the main text.

\FloatBarrier
\section{Shock-definition and horizon robustness}

The shock definitions anchor to a 1990--2000 pre-study climatology: temperature and precipitation z-scores use baseline means and standard deviations from that window, and the heat and cold burden indicators use national 80th-percentile degree-day cutoffs computed over it. Extreme drought is defined by an absolute threshold - a Palmer index at or below -4 - and does not depend on the baseline. To assess sensitivity to the baseline horizon, I rebuild the z-scores and degree-day cutoffs under baselines of 1990--2005 and 1990--2010, both of which end before the 2011--2023 outcome window and so introduce no look-ahead, and re-estimate the affected headline specifications; the 1990--2000 rebuild first reproduces the shipped shock flags exactly (zero mismatches in 40,781 county-years) and the published coefficients to the digit. The definitions prove insensitive to the choice. The national cutoffs move by about two percent at most, the share of county-years flagged shifts by at most 1.1 percentage points, the Medicare heat responses are stable (\$112, \$107, and \$113 per beneficiary contemporaneous and \$176, \$179, and \$180 at one lag across the three baselines, with emergency-visit responses similar), and the county cold-employment response is likewise stable (-714, -607, and -613 employed persons at the two-year lag; p = 0.035 to 0.039). The one quantifiable sensitivity is the state cold--debt coefficient, which attenuates from 1.35 percentage points (p = 0.012) under the 1990--2000 baseline to 1.03 (p = 0.027) and 0.85 (p = 0.044) as the baseline lengthens - sign-stable and significant at the five-percent level under every choice, but smaller, because a baseline that includes the warming 2000s raises the reference mean, and the z-score threshold then selects milder cold years. I therefore cite that response as a 0.85 to 1.35 percentage-point range wherever baseline robustness is at issue.

The dynamic specifications are similarly insensitive to the estimation horizon. Re-running the event-study machinery with maximum horizons of two through five years leaves every headline coefficient at horizons zero through two within one standard error of its reported value, with no sign changes; the only consequential choice is truncation to a two-year horizon, which understates the cold-employment response. The longer horizons add information rather than fragility: the drought-debt response is transient by a four-year horizon, while the cold-employment response persists at three and four years (p = 0.009 and 0.006) - dynamics taken up in the second essay.

\FloatBarrier
\clearpage
\bibliographystyle{apalike}
\bibliography{../references}

\end{document}